\section{The cut-offs in front of the Gaussian time average}
\label{sec:hastings_moved_estimate}

% Hastings, arXiv:0705.2024v4, law16.tex, lines 865--876 (eq. moved), using
% eq. obd (lines 836--840). Finite complex matrices; no lattice content.

After the Gaussian time average has been moved from the Hamiltonian to the
nearby operator $M=M_L+M_B+M_R$, the proof of Lemma~1 of
\cite{Hastings2007AreaLaw} multiplies it on the right by the spectral cut-offs
$O_L$ and $O_R$. Since the cut-offs nearly fix the ground state, this changes
the distance to the ground-state projector only slightly. This section states
the estimate for finite complex matrices with explicit constants. Norms are
operator norms.

\begin{lemma}[A nearly invariant projection next to the ground-state projector]
    \label{lem:hastings_ground_projector_mul}
    \lean{TNLean.MPS.AreaLaw.Hastings.norm_groundProjector_mul_sub_le}
    \leanok
    Let $O$ be a self-adjoint matrix and $\Omega$ a unit vector. Then
    \begin{align*}
        \bigl\||\Omega\rangle\langle\Omega|\,O-|\Omega\rangle\langle\Omega|\bigr\|
        \le\|\Omega-O\Omega\|.
    \end{align*}
\end{lemma}
\begin{proof}\leanok
    Taking adjoints, the norm equals that of
    $(O-1)|\Omega\rangle\langle\Omega|$, which maps $x$ to
    $\langle\Omega,x\rangle(O\Omega-\Omega)$. The Cauchy--Schwarz inequality
    bounds its norm by $\|x\|\,\|O\Omega-\Omega\|$.
\end{proof}

\begin{theorem}[The cut-offs next to the Gaussian time average]
    \label{thm:hastings_moved_estimate}
    \lean{TNLean.MPS.AreaLaw.Hastings.norm_gaussianTimeAverage_mul_sub_groundProjector_le,
        TNLean.MPS.AreaLaw.Hastings.norm_gaussianTimeAverage_mul_mul_sub_groundProjector_le}
    \leanok
    \uses{def:hastings_gaussian_time_average,
        thm:hastings_gaussian_time_average_gap}
    Let $H$, $\Omega$ and $\Delta E>0$ satisfy the hypotheses of
    Theorem~\ref{thm:hastings_gaussian_time_average_gap}, and let $M$ be
    Hermitian with $\|M-H\|\le\delta$. Let $O_L$ and $O_R$ be commuting
    orthogonal projections with $\|\Omega-O_LO_R\Omega\|\le\eta$. Then for
    every real $q$,
    \begin{align*}
        \bigl\|P_q(M)\,O_LO_R-|\Omega\rangle\langle\Omega|\bigr\|
        \le e^{-q/2}+\frac{\sqrt q}{\Delta E}\,\delta+\eta.
    \end{align*}
    The same bound holds with $O_LO_R$ replaced by a single orthogonal
    projection. With $M=M_L+M_B+M_R$, $\delta$ of order
    $\Delta E^{-1}J^2e^{-l/3\xi}$, and $\eta$ of order $e^{-l/6\xi}$ by the
    estimate on the spectral cut-offs, this is the estimate on
    \cite[lines 865--876]{Hastings2007AreaLaw}; the term $e^{-q/2}$ is the gap
    bound of \cite[line 848]{Hastings2007AreaLaw}.
\end{theorem}
\begin{proof}\leanok
    \uses{thm:hastings_gaussian_time_average_stability,
        lem:hastings_ground_projector_mul}
    Write $P_0=|\Omega\rangle\langle\Omega|$ and $O=O_LO_R$, an orthogonal
    projection since the factors commute. Then
    $P_q(M)O-P_0=(P_q(M)-P_0)O+(P_0O-P_0)$. Since $\|O\|\le1$,
    Theorem~\ref{thm:hastings_gaussian_time_average_stability} bounds the
    first term by $e^{-q/2}+\sqrt q\,\delta/\Delta E$, and
    Lemma~\ref{lem:hastings_ground_projector_mul} bounds the second by
    $\eta$.
\end{proof}
